% ===== Джерела до розділу 8. Ключі з префіксом ch08: =====
% Позиції ядра (bib-core.tex) тут НЕ дублюються.

% --- Пакування та розпакування ---
\bibitem{ch08:ugarte} Ugarte-Pedrero X., Balzarotti D., Santos I., Bringas P. G. SoK: Deep Packer Inspection --- A Longitudinal Study of the Complexity of Run-Time Packers. \emph{2015 IEEE Symposium on Security and Privacy}. San Jose: IEEE, 2015. P.~659--673. DOI: 10.1109/SP.2015.46.

\bibitem{ch08:guo} Guo F., Ferrie P., Chiueh T. A Study of the Packer Problem and Its Solutions. \emph{Recent Advances in Intrusion Detection (RAID 2008)}. LNCS Vol.~5230. Berlin: Springer, 2008. P.~98--115.

\bibitem{ch08:royal} Royal P., Halpin M., Dagon D., Edmonds R., Lee W. PolyUnpack: Automating the Hidden-Code Extraction of Unpack-Executing Malware. \emph{22nd Annual Computer Security Applications Conference (ACSAC)}. Miami Beach: IEEE, 2006. P.~289--300.

% --- Обфускація ---
\bibitem{ch08:collberg} Collberg C., Nagra J. Surreptitious Software: Obfuscation, Watermarking, and Tamperproofing for Software Protection. Boston: Addison-Wesley, 2009. 792~p.

\bibitem{ch08:oreans} Schrittwieser E., Katzenbeisser S., Kinder J., Merzdovnik G., Weippl E. Protecting Software through Obfuscation: Can It Keep Pace with Progress in Code Analysis? \emph{ACM Computing Surveys}. 2016. Vol.~49, №~1. Art.~4. P.~1--37.

\bibitem{ch08:junod} Junod P., Rinaldini J., Wehrli J., Michielin J. Obfuscator-LLVM --- Software Protection for the Masses. \emph{2015 IEEE/ACM 1st International Workshop on Software Protection (SPRO)}. Florence: IEEE, 2015. P.~3--9.

% --- Поліморфізм та метаморфізм ---
\bibitem{ch08:szor} Szor P. The Art of Computer Virus Research and Defense. Boston: Addison-Wesley, 2005. 744~p.

\bibitem{ch08:you} You I., Yim K. Malware Obfuscation Techniques: A Brief Survey. \emph{2010 International Conference on Broadband, Wireless Computing, Communication and Applications}. Fukuoka: IEEE, 2010. P.~297--300.

% --- Антианаліз ---
\bibitem{ch08:branco} Branco R. R., Barbosa G. N., Neto P. D. Scientific but Not Academical Overview of Malware Anti-Debugging, Anti-Disassembly and Anti-VM Technologies. \emph{Black Hat USA}. Las Vegas, 2012.

% --- Деобфускація ---
\bibitem{ch08:yadegari} Yadegari B., Johannesmeyer B., Whitely B., Debray S. A Generic Approach to Automatic Deobfuscation of Executable Code. \emph{2015 IEEE Symposium on Security and Privacy}. San Jose: IEEE, 2015. P.~674--691.

% --- Українські джерела ---
\bibitem{ch08:obfusc-ua} Методи обфускації програмного коду та підходи до їх подолання в задачах аналізу шкідливого програмного забезпечення. \emph{Кібербезпека: освіта, наука, техніка}. 2024. URL: \url{https://csecurity.kubg.edu.ua/index.php/journal/article/view/702}.

\bibitem{ch08:packer-ua} Виявлення запакованого шкідливого програмного забезпечення на основі аналізу ентропії та структурних ознак. \emph{Захист інформації}. 2023. Т.~25, №~3.
